\section*{Chapitre \ref{ch:nw:long-haul-fibre} --- Fibre longue distance}

\begin{solution}{exo:nw:long-haul-fibre:1}
$0.16 \times 80 = \qty{12.8}{\decibel}$, soit un facteur $10^{-1.28}$~: environ 5.2~\% de la lumière subsiste.
\end{solution}

\begin{solution}{exo:nw:long-haul-fibre:2}
$(6\,550 - 21.7) \times 10^{-6} \times c_0 / 1.462 = \qty{1339}{\kilo\metre}$ de verre, un \omterm{def:nw:the-north-american-map:factor}{facteur de tracé} de 1.185 par rapport à la
\omterm{def:nw:the-north-american-map:geodesic}{distance géodésique} de \qty{1129.2}{\kilo\metre}.
\end{solution}

\begin{solution}{exo:nw:long-haul-fibre:3}
Avec une \omterm{def:nw:long-haul-fibre:dark}{fibre noire}, l'acheteur choisit et contrôle les équipements (transpondeurs, \omterm{def:nw:long-haul-fibre:fec}{correction d'erreurs directe}, \omterm{def:nw:long-haul-fibre:disp}{compensation de dispersion}, sans commutateurs), donc il contrôle chaque
microseconde au-delà du verre~; avec un \omterm{def:nw:long-haul-fibre:dark}{service de longueur d'onde}, l'équipement du vendeur et ses choix fixent la latence, et peuvent changer lorsque le vendeur
le modernise.
\end{solution}

\begin{solution}{exo:nw:long-haul-fibre:4}
La compensation optique ajoute une longueur de fibre compensatrice de dispersion à chaque tronçon, et la lumière met du temps à la traverser~: dans le modèle, un cinquième du
trajet, soit plus d'une milliseconde sur le tracé de Chicago. La compensation électronique annule la dispersion par calcul dans le récepteur, ce qui coûte un
délai de traitement à l'extrémité du tracé, indépendant de sa longueur.
\end{solution}

\begin{solution}{exo:nw:long-haul-fibre:5}
À travers l'océan, un câble peut suivre de près la \omterm{def:nw:the-north-american-map:geodesic}{distance géodésique}, avec des détours seulement aux stations d'atterrissement et aux liaisons vers les centres de données~; sur terre, le tracé
doit suivre les droits de passage qu'il peut obtenir (routes, voies ferrées, couloirs de services publics), en contournant le relief et les propriétés. Le tracé atlantique est aussi
cinq fois plus long, si bien que les détours fixes à ses extrémités pèsent moins en proportion.
\end{solution}

\begin{solution}{exo:nw:long-haul-fibre:6}
Si les deux sens empruntent les mêmes fibres, le même tracé et les mêmes équipements (sinon les deux moitiés diffèrent, et un aller-retour ne peut pas les distinguer), si la
mesure va de baie à baie ou entre les équipements du vendeur, et s'il s'agit d'un chiffre typique, optimal ou garanti.
\end{solution}

\begin{solution}{exo:nw:long-haul-fibre:7}
Le tracé hertzien à 1.011 dans l'air met $1.011 \times 3\,767.8 = \qty{3809}{\micro\second}$ entre les deux bâtiments. La \omterm{def:nw:long-haul-fibre:hcf}{fibre à cœur creux} suivant
la \omterm{def:nw:the-north-american-map:geodesic}{distance géodésique} met \qty{3778}{\micro\second}~; avec les \qty{21.7}{\micro\second} d'équipement du modèle, le \omterm{def:nw:the-north-american-map:factor}{facteur de tracé} de la \omterm{def:nw:long-haul-fibre:hcf}{fibre à cœur creux} doit être au plus
$(3\,809 - 21.7)/3\,778 = 1.0025$ (1.008 sans équipement). Une \omterm{def:nw:long-haul-fibre:hcf}{fibre à cœur creux} devrait suivre un tracé aussi droit que la liaison hertzienne et
n'ajouter presque aucun délai d'équipement.
\end{solution}

\begin{solution}{exo:nw:long-haul-fibre:8}
Un tracé plus long de 3~\% que la \omterm{def:nw:the-north-american-map:geodesic}{distance géodésique} dans le verre met $1.03 \times 1.462 = 1.51$ fois le plancher du vide~; une liaison hertzienne avec un \omterm{def:nw:the-north-american-map:factor}{facteur de tracé} de 1.01
à 1.10 dans l'air met de 1.01 à 1.10 fois ce plancher. Le zigzag des pylônes coûte quelques pour cent~; le verre coûte 46~\%. La rectitude ne peut pas
compenser le milieu.
\end{solution}

\begin{solution}{pb:nw:long-haul-fibre:1}
\textbf{Partie I.}
\begin{enumerate}
\item \qty{1129.2}{\kilo\metre}~; \qty{3.767}{\milli\second} dans le vide, \qty{5.507}{\milli\second} dans la fibre.
\item \qty{13.3}{\milli\second} d'aller-retour, contre un plancher d'aller-retour en fibre de \qty{11.014}{\milli\second}.
\item $6\,650 - 5\,507 = \qty{1143}{\micro\second}$.
\item Parce que le tracé est en verre~: même le long de la \omterm{def:nw:the-north-american-map:geodesic}{distance géodésique}, il ne peut pas battre $n_g d/c_0$~; le plancher du vide mesure ce que coûte le milieu, et non
  ce que le tracé gaspille.
\end{enumerate}
\textbf{Partie II.}
\begin{enumerate}[resume]
\item \qty{1359}{\kilo\metre} de verre, en 17 tronçons d'au plus \qty{80}{\kilo\metre}.
\item \qty{21.7}{\micro\second} (deux terminaux et 17 amplificateurs)~: 1.9~\% de l'excès.
\item \qty{1121}{\micro\second} de trajet supplémentaire, \qty{230}{\kilo\metre} de verre.
\item Le trajet se réduit à \qty{1323}{\kilo\metre}, soit encore \qty{193}{\kilo\metre} (\qty{943}{\micro\second}) de plus que la \omterm{def:nw:the-north-american-map:geodesic}{distance géodésique}~: la plus grande part
  de l'excès reste dans le trajet.
\end{enumerate}
\textbf{Partie III.}
\begin{enumerate}[resume]
\item \qty{1.33}{\milli\second}~: \qty{7.976}{\milli\second} au total.
\item \qty{4.569}{\milli\second}, soit \qty{2.08}{\milli\second} de moins.
\item \qty{5.528}{\milli\second}~: le redressement du tracé ferait gagner \qty{1.12}{\milli\second}, la \omterm{def:nw:long-haul-fibre:hcf}{fibre à cœur creux} \qty{2.08}{\milli\second}.
\item La liaison hertzienne rejoint Carteret depuis Aurora, \qty{52}{\kilo\metre} plus à l'ouest, en \qty{3.982}{\milli\second}~: \qty{0.59}{\milli\second}
  plus rapide que la \omterm{def:nw:long-haul-fibre:hcf}{fibre à cœur creux} sur le tracé de 2010, \qty{2.67}{\milli\second} plus rapide que le tracé de 2010.
\end{enumerate}
\textbf{Partie IV.}
\begin{enumerate}[resume]
\item \emph{Résultat nommé}~: sur l'aller-retour de \qty{13.3}{\milli\second}, \qty{11.01}{\milli\second} sont du verre le long de la \omterm{def:nw:the-north-american-map:geodesic}{distance géodésique},
  \qty{2.24}{\milli\second} de trajet supplémentaire et \qty{0.04}{\milli\second} d'équipement (avec les hypothèses du modèle)~; une \omterm{def:nw:long-haul-fibre:hcf}{fibre à cœur creux} sur le même trajet ferait
  gagner \qty{2.08}{\milli\second} à l'aller, \qty{4.16}{\milli\second} en aller-retour.
\item La répartition entre trajet et équipement~: chaque microseconde d'équipement supposée représente \qty{205}{\metre} de trajet en moins.
\item Les réseaux hertziens suivant presque la même \omterm{def:nw:the-north-american-map:geodesic}{distance géodésique}, dans un milieu 46~\% plus rapide, sont restés à 1 à 2~\% du plancher de l'air.
\item Bande passante et fiabilité~: elle transporte les données de marché et les ordres par tous les temps et sert de secours aux liaisons radio.
\item Cela dépend de la durée de vie de la stratégie et du contrôle~: pour une société qui doit contrôler chaque équipement pendant des années, la \omterm{def:nw:long-haul-fibre:dark}{fibre noire}~; pour une
  sauvegarde ou un tracé de données, une longueur d'onde louée.
\item Un \omterm{def:nw:the-north-american-map:factor}{facteur de tracé} à moins d'un pour cent de la \omterm{def:nw:the-north-american-map:geodesic}{distance géodésique} et presque aucun délai d'équipement (voir l'exercice~\ref{exo:nw:long-haul-fibre:7}).
\item C'est sur terre que le temps se perd dans le trajet~: le câble atlantique, qui n'a pas de droits de passage à suivre, a un facteur de 1.09, contre
  1.20 pour le tracé de Chicago.
\item Onze vont à la lumière dans le verre le long de la \omterm{def:nw:the-north-american-map:geodesic}{distance géodésique}, deux aux détours du trajet, et presque rien à l'équipement.
\end{enumerate}
\end{solution}

\begin{solution}{iq:nw:long-haul-fibre:1}
À $c_0/1.46$, soit environ \qty{205000}{\kilo\metre\per\second}~: \qty{4.88}{\micro\second} par kilomètre, \qty{4.88}{\milli\second} par
\qty{1000}{\kilo\metre} à l'aller, contre \qty{3.34}{\milli\second} dans le vide ou dans l'air.
\iqlookfor{L'\omterm{def:nw:the-north-american-map:floor}{indice de groupe}~; aller simple contre aller-retour~; la règle empirique par kilomètre.}
\end{solution}

\begin{solution}{iq:nw:long-haul-fibre:2}
Les détours du trajet par rapport à la \omterm{def:nw:the-north-american-map:geodesic}{distance géodésique}, la \omterm{def:nw:long-haul-fibre:disp}{compensation de dispersion} optique (du verre en plus), les amplificateurs (faibles), les terminaux~: transpondeurs et
\omterm{def:nw:long-haul-fibre:fec}{correction d'erreurs directe}, et les éventuels commutateurs ou routeurs en chemin~; plus les raccordements entre les extrémités du tracé et les baies.
\iqlookfor{La longueur du trajet domine~; DCF contre compensation électronique~; la FEC comme compromis entre distance et délai.}
\end{solution}

\begin{solution}{iq:nw:long-haul-fibre:3}
Les ondes hertziennes se propagent dans l'air à presque $c_0$, 46~\% plus vite que la lumière dans le verre, le long d'une ligne presque droite de pylônes~; la fibre suit les droits de passage
et elle est plus lente au kilomètre. La fibre reste utile pour la bande passante, la fiabilité par mauvais temps, et tout ce qui n'est pas
les quelques messages sensibles à la latence.
\iqlookfor{Milieu et trajet~; capacité et météo~; les deux utilisés ensemble.}
\end{solution}

\begin{solution}{iq:nw:long-haul-fibre:4}
Un droit d'usage exclusif à long terme, généralement payé d'avance, sur des fibres déterminées ou une capacité, sans posséder le câble. Une société achète une \omterm{def:nw:long-haul-fibre:dark}{fibre noire} lorsqu'elle a besoin
de contrôler les équipements et donc la latence, prévoit d'utiliser le tracé plus longtemps que la durée de seuil de rentabilité, et dispose des équipes pour l'exploiter~;
sinon elle loue.
\iqlookfor{Contrôle de la latence~; coût initial contre coût mensuel~; obligations de maintenance~; durée de vie de la stratégie.}
\end{solution}

\begin{solution}{iq:nw:long-haul-fibre:5}
Fibre qui guide la lumière dans un cœur d'air, si bien que la lumière se propage à presque $c_0$~; les conceptions récentes perdent aussi moins de lumière que le verre plein. Pour le trading, elle
réduirait d'environ 30~\% la latence de la fibre, rapprochant les tracés en fibre du plancher de l'air avec la bande passante de la fibre et son indépendance à la météo, et
menaçant l'avantage des liaisons hertziennes sur les tracés qui peuvent être construits en ligne droite.
\iqlookfor{La latence et les pertes à la fois~; le tracé doit toujours être droit~; déploiement et coût comme questions ouvertes.}
\end{solution}

\begin{solution}{iq:nw:long-haul-fibre:6}
Les extrémités diffèrent de celles qui ont été calculées (un bâtiment plus proche)~; le chiffre n'est pas un aller simple (un aller-retour divisé par deux sur des trajets asymétriques, ou
une erreur d'unité)~; le tracé n'est pas entièrement en fibre standard (\omterm{def:nw:long-haul-fibre:hcf}{fibre à cœur creux}, ou partie radio)~; ou la mesure est fausse (horloges non synchronisées). Il faut demander
les extrémités, la méthode et le milieu, puis mesurer.
\iqlookfor{La physique comme contrôle de cohérence~; extrémités et définitions~; mesure.}
\end{solution}
